\documentclass[border=10pt,tikz]{standalone}
\usepackage{amsmath} % for \text
\usepackage{tikz}
\usepackage{tikz-3dplot}
\usetikzlibrary{math}
\usepackage[active,tightpage]{preview}
\PreviewEnvironment{tikzpicture}
\setlength\PreviewBorder{1pt}
%
% File name: surface-of-revolution.tex
% Description: 
% A geometric representation of a surface of revolution is shown.
% 
% Date of creation: October, 17th, 2021.
% Date of last modification: October, 9th, 2022.
% Author: Efraín Soto Apolinar.
% https://www.aprendematematicas.org.mx/author/efrain-soto-apolinar/instructing-courses/
% Source: page 450 of the 
% Glosario Ilustrado de Matem\'aticas Escolares.
% https://tinyurl.com/5udm2ufy
%
% Terms of use:
% According to TikZ.net
% https://creativecommons.org/licenses/by-nc-sa/4.0/
% Your commitment to the terms of use is greatly appreciated.
%
\begin{document}
	%
	\tdplotsetmaincoords{70}{30}
	\begin{tikzpicture}[tdplot_main_coords]
		% The function that is rotated
		\tikzmath{function f(\x) {return 1.5 - 0.35*sin(\x r);};}
        \def\R{3.0}
        \def\th{20}
		% \pgfmathsetmacro{\dominio}{2.0}
		% \pgfmathsetmacro{\xi}{-\dominio}
		% \pgfmathsetmacro{\step}{(\dominio-\xi)/70.0}
		% \pgfmathsetmacro{\xs}{\xi+\step}
		% \pgfmathsetmacro{\max}{3}
		% Circumferences (behind the coordiante axis)
		% \foreach \x in {\xi,\xs,...,\dominio}{
		% 	\pgfmathsetmacro{\radio}{f(\x)}	% radius of the circumference of the solid of revolution
		% 	\draw[cyan,very thick,opacity=0.35] plot[domain=0.5*pi:2.0*pi,smooth,variable=\t] ({\radio*cos(\t r)},{\x},{\radio*sin(\t r});	
		% }
		% % Part of the solid of revolution behind the coordinate axis
		% \foreach \angulo in {358,356,...,90}{
		% 	\draw[cyan,very thick,rotate around y=\angulo,opacity=0.35] plot[domain=-\dominio:\dominio,smooth,variable=\t] ({0},{\t},{f(\t)});
		% }
		% % Graph of the function rotated about the $y$ axis

        \node [] at (-0.4, 0, 3) {(a)};

        \draw[very thick] (0, 0, -3) -- (0, 0, 3) node [right] {$O$};
 

        \draw[thick] ({\R*cos(-\th)}, {\R*sin(-\th)}, 2) arc (-\th:{\th}:\R);

        \draw[red, thick, <->] ({\R*cos(-\th)/2}, {\R*sin(-\th)/2}, 2) arc (-\th:{\th}:{\R/2}) node [midway, right, yshift=-3] {$\alpha$};


        \draw[thick] ({\R*cos(-\th)}, {\R*sin(-\th)}, -2) arc (-\th:{\th}:\R);

  
        \draw[thick] (0, 0, 2) -- ({\R*cos(-\th)}, {\R*sin(-\th)}, 2) 
            .. controls ({\R*cos(-\th) *1.2}, {\R*sin(-\th) * 1.2}, 1) and ({\R*cos(-\th) * 1.2}, {\R*sin(-\th) * 1.2}, -1) 
            .. ({\R*cos(-\th)}, {\R*sin(-\th)}, -2)
            -- (0, 0, -2) -- cycle;

        \draw[thick] (0, 0, 2) -- ({\R*cos(\th)}, {\R*sin(\th)}, 2) 
            .. controls ({\R*cos(\th) *1.2}, {\R*sin(\th) * 1.2}, 1) and ({\R*cos(\th) * 1.2}, {\R*sin(\th) * 1.2}, -1) 
            .. ({\R*cos(\th)}, {\R*sin(\th)}, -2);
            
        \node [] at ({\R*cos(\th)*1.2}, {\R*sin(\th)*1.2}, 1) {C};

        \draw[thick, dashed] (0, 0, 0) -- ({\R*cos(-\th)*1.1}, {\R*sin(-\th)*1.1}, 0) arc (-\th:{\th}:{\R*1.1}) --cycle node[midway, above] {$r$};
            


        \begin{scope}[xshift=150]

        \node [] at (-0.4, 0, 3) {(b)};

        \draw[very thick] (0, 0, -3) -- (0, 0, 3) node [right] {$O$};
 

        \draw[thick] ({\R*cos(-\th)}, {\R*sin(-\th)}, 2) arc (-\th:{\th}:\R);

        \draw[red, thick, <->] ({\R*cos(-\th)/2}, {\R*sin(-\th)/2}, 2) arc (-\th:{\th}:{\R/2}) node [midway, right, yshift=-3] {$\alpha$};


        \draw[thick] ({\R*cos(-\th)}, {\R*sin(-\th)}, -2) arc (-\th:{\th}:\R);

  
        \draw[thick, fill=gray!30] (0, 0, 2) -- ({\R*cos(-\th)}, {\R*sin(-\th)}, 2) 
            .. controls ({\R*cos(-\th) *1.2}, {\R*sin(-\th) * 1.2}, 1) and ({\R*cos(-\th) * 1.2}, {\R*sin(-\th) * 1.2}, -1) 
            .. ({\R*cos(-\th)}, {\R*sin(-\th)}, -2)
            -- (0, 0, -2) -- cycle;

        \draw[thick] (0, 0, 2) -- ({\R*cos(\th)}, {\R*sin(\th)}, 2) 
            .. controls ({\R*cos(\th) *1.2}, {\R*sin(\th) * 1.2}, 1) and ({\R*cos(\th) * 1.2}, {\R*sin(\th) * 1.2}, -1) 
            .. ({\R*cos(\th)}, {\R*sin(\th)}, -2);
            
        \node [] at ({\R*cos(\th)*1.2}, {\R*sin(\th)*1.2}, 1) {C};

        \draw[thick, dashed] (0, 0, 0) -- ({\R*cos(-\th)*1.1}, {\R*sin(-\th)*1.1}, 0) arc (-\th:{\th}:{\R*1.1}) --cycle node[midway, above] {$r$};


        \end{scope}
		% \draw[red,ultra thick] plot[domain=-\dominio:\dominio,smooth,variable=\t] ({0},{\t},{f(\t)}) node [above right] {$z = f(y)$};
		% % Coordinate axis
		% \draw[thick,->] (0,0,0) -- (0,\max,0) node [right] {$y$};
		% \draw[thick,->] (0,0,0) -- (\max,0,0) node [left] {$x$};
		% \draw[thick,->] (0,0,0) -- (0,0,\max) node [above] {$z$};
		% % Circumferences (in front of the coordiante axis)
		% \foreach \x in {\xi,\xs,...,\dominio}{
		% 	\pgfmathsetmacro{\radio}{f(\x)}	% Radio del círculo al inicio del sólido de revolución
		% 	\draw[cyan,very thick,opacity=0.35] plot[domain=0.0:0.5*pi,smooth,variable=\t] ({\radio*cos(\t r)},{\x},{\radio*sin(\t r});	
		% }
		% % The solid of revolution (in front of the coordinate axis)
		% \foreach \angulo in {0,2,...,89}{
		% 	\draw[cyan,very thick,rotate around y=\angulo,opacity=0.35] plot[domain=-\dominio:\dominio,smooth,variable=\t] ({0},{\t},{f(\t)});
		% }
	\end{tikzpicture}
\end{document}